% Page 049 — pages imprimées 45–46
% Contenu seul : le préambule, l'en-tête et la pagination sont gérés par meyrueis.tex

\textbf{Climat général en 1932 -- 1935}

De profonds changements se dessinèrent au cours de ces trois années. Les mutations
ne datent pas d’aujourd’hui seulement... la radio s’étendait partout et prenait une
grande place dans la vie de nos campagnes. Nous avions acquis un petit poste « Point
Bleu » bien insuffisant et modeste, mais nous l’écoutions tous les jours. Grâce à une
antenne tendue depuis le clocher du temple jusqu’au toit du presbytère, en l’absence
de moteurs et d’appareils ménagers, je captais les émissions avec une grande pureté.
C’était l’époque où les vedettes se nommaient Gaby Morlay, Sim Viva ou la
danseuse Mitsi, beautés lointaines et inconnues mais derrière le nom desquelles
j’imaginais la grâce féminine parisienne. L’époque également, hélas, des grands
scandales financiers, Stavisky et autres.

Puis ce fut le 6 Février 1934 et surtout le début des hurlements sinistres d’un inconnu
qui se nommait Adolf Hitler. Vraiment j’avais le sentiment que la relative tranquillité
que nous avions connue depuis l’armistice de 1919 était, bel et bien, en train de
disparaître.

Voilà un éclairage bien incomplet de ce temps de Meyrueis. Je pourrais développer.
Parler de Millau, de Florac, de nos veillées à Gatuzières ou à Cabrillac, de nos
campagnes « Christianisme Social », de nos représentations théâtrales avec les jeunes
etc Tout cela fut fort intéressant, remarquable même. Mais pourquoi le temps passé à
Meyrueis fut-il si court ?...

\emph{[... ( Jean Gounelle, comme plusieurs membres de sa famille, était attiré par l’Afrique du Nord :
l’énigme de ces pays évangélisés aux premiers siècles, leur essor, puis l’élimination à peu près totale
des églises chrétiennes par l’Islam du 8\textsuperscript{o} au 12\textsuperscript{o} siècle... Ne serait-ce pas le moment de s’y
réimplanter ? Oran n’avait plus de pasteur... Il fut nommé... et partit sans attendre... )....]}

Mais, cinquante ans plus tard, je conserve au fond de mon cœur, empreint d’une
certaine nostalgie, l’apport précieux que fut pour nous notre bref séjour à Meyrueis.
Et quand j’y retourne maintenant de Mandagout, (\emph{près du Vigan}) je perçois mieux, me
semble-t-il, les richesses et les bénédictions reçues dans cette humble cité ! Je m’en
souviens comme on se souvient avec émotion de ces périodes de préparation intense,
aussi bien intérieure que sur le terrain, où, loin des apparences et guidé par une main
invisible, se joue vraiment le destin d’un homme.
